\section{Experiments}\label{sec:experiments}
We evaluate the pipeline on three synthetic datasets, each designed to mimic the statistical properties of a distinct domain while keeping the computational load low enough to run on a standard laptop within seconds.
This allows a fully reproducible proof‑of‑concept; the same pipeline can be applied to real‑world data without modification.

\subsection{ECG: Normal vs. Arrhythmic}
We generate a two‑class synthetic ECG dataset consisting of a sinusoidal signal with additive noise (normal) and a pure noise signal (arrhythmic).
Both segments contain 10,000 samples; windows of length $W=360$ are extracted with stride $S=180$, yielding 107 windows.
The binary target labels are $0$ (normal) and $1$ (arrhythmic).

\subsection{Financial: Sector‑like random walks}
We create 33 synthetic stock price series by cumulating Gaussian random walks with small, sector‑dependent drifts.
Windows of length $W=20$ with stride $S=5$ are extracted from the concatenated series, resulting in 165 windows.
Each window is labelled by the sector index (0‑10) of the stock it belongs to, mimicking a sector‑classification task.

\subsection{Seismic: Noise vs. P‑wave event}
A synthetic seismic trace is constructed by concatenating pure noise, an event segment containing a damped sinusoidal P‑wave pulse, and further noise.
Each segment contains 6,000 samples; windows of $W=600$ with stride $S=300$ are extracted, yielding 38 windows.
Labels are $0$ (noise) and $1$ (event), and the dataset is balanced by truncation.

For all datasets, we compute the adjacency matrix as described in Section~\ref{sec:method} and run Louvain clustering with default resolution.
An ablation study with 3 independent train‑test splits (70/30) is performed to compute the Functional Significance Scores $S_k$.